\subsection{帰宅予測時刻の取得}
% 帰宅推奨時刻および提示時刻の算出に用いる，各滞在者の帰宅予測時刻の取得方法について述べる．

本システムでは，先行研究で提案されているシステム\cite{hayashi}が提供する個人単位の帰宅予測時刻を取得し，入力として利用する．
% このシステムではユーザの滞在履歴から来訪時刻および帰宅時刻のみを抽出し，これらに対して GMM（Gaussian Mixture Model）によるクラスタリングを行う．
% 得られた各成分の累積分布関数を対応するデータ数で重み付けして合成し，一定の確率を超える時刻を帰宅予測時刻として算出しており，これをWeb APIを通じて提供している．
現在滞在している各ユーザについて，このWeb APIから帰宅予測時刻を取得し，複数人の帰宅行動を集約するための基礎データとして用いる．
先行研究により，実測値に近い精度で未来の在室情報を予測できていると示されているため，本研究ではその結果を活用し複数人が同時に帰宅しやすい時刻を導出するための基礎情報として位置付けている．
すなわち，本研究の焦点は帰宅予測そのものではなく，妥当な個人の在室予測を前提として，複数人が同時に帰宅できる可能性の高い時刻を導出する点にある．


時間の経過に伴い，滞在者の人数や構成，ならびに各ユーザの帰宅予測時刻は変化する．
そのため，本システムでは一定間隔で帰宅予測時刻を再取得し，最新の状況を反映した帰宅推奨時刻および提示時刻を算出する．
再取得の間隔は，帰宅行動の変化を過度に遅延なく反映しつつ，不要な再計算を避ける観点から30分と設定した．
また，帰宅推奨時刻が近づいたユーザに対しては，十分な意思決定時間を確保した上で提示が行われるよう，更新間隔と判定条件を設計している
なお，更新間隔および判定条件は利用環境や運用方針に応じて調整可能なパラメータであり，本研究では対象環境における運用を想定して30分を採用した．